% Tufte 的可选参数包装隐藏了脚注正文；为 biblatex 提供可探测的颜色分组。
% 委托原实现，保留原有边注定位和可选垂直偏移接口。
\makeatletter
\let\book@tufte@footnotetext\@footnotetext
\def\@footnotetext{\color@begingroup\book@footnote@dispatch}
\newcommand{\book@footnote@dispatch}[2][0pt]{%
  \book@tufte@footnotetext[#1]{#2}\color@endgroup}
\makeatother
\usepackage[backend=biber,style=authortitle,giveninits=true,
  maxnames=99,minnames=99,uniquename=false,doi=false,url=false,isbn=false]{biblatex}
\addbibresource{tex/references.bib}
\DeclareNameAlias{author}{given-family}
\DeclareNameAlias{editor}{given-family}
\DeclareDelimFormat{multinamedelim}{\addcomma\space}
\DeclareDelimFormat{finalnamedelim}{\addcomma\space}
\DeclareFieldFormat[book,article,inproceedings]{title}{\mkbibemph{#1}}
\DeclareFieldFormat[article]{journaltitle}{#1}
\DeclareFieldFormat[inproceedings]{booktitle}{#1}
\DeclareFieldFormat[book]{edition}{\mkbibordedition{#1}\addspace ed.}
\DeclareFieldFormat[article]{volume}{vol.\addspace #1}
\DeclareFieldFormat[article]{number}{no.\addspace #1}
\DeclareFieldFormat[article,inproceedings]{pages}{pp.\addspace #1}
\renewcommand{\newunitpunct}{\addcomma\space}
\newbibmacro*{book:author-title}{%
  \ifnameundef{author}{%
    \ifnameundef{editor}{}{%
      \printnames{editor}\setunit{\addcomma\space}\printtext{eds.}}%
  }{\printnames{author}}%
  \newunit\printfield{title}}
\DeclareBibliographyDriver{book}{%
  \usebibmacro{book:author-title}\newunit\printlist{publisher}
  \newunit\printfield{year}\newunit\printfield{edition}\usebibmacro{finentry}}
\DeclareBibliographyDriver{article}{%
  \usebibmacro{book:author-title}\setunit{\addspace}
  \printtext{In\addspace}\printfield{journaltitle}
  \newunit\printfield{year}\newunit\printfield{volume}
  \newunit\printfield{number}\newunit
  \iffieldundef{pages}{\printfield{eid}}{\printfield{pages}}\usebibmacro{finentry}}
\DeclareBibliographyDriver{inproceedings}{%
  \usebibmacro{book:author-title}\setunit{\addspace}
  \printtext{In\addspace}\printfield{booktitle}
  \newunit\printfield{year}\newunit\printfield{pages}\usebibmacro{finentry}}
% 无排序压缩；每篇文献使用独立编号和边注，可带前注和定位信息。
% \sidecite[参见][第 3 章]{bishop2006}；\cite 使用完全相同的接口。
\makeatletter
\ExplSyntaxOn
\cs_new_protected:Npn \book_sidecite_note:n #1
  {
    \stepcounter\@mpfn
    \protected@xdef\@thefnmark{\thempfn}
    \@footnotemark
    \book@tufte@footnotetext[0pt]
      {\color@begingroup\booknotefont #1\color@endgroup}
  }
\cs_new_protected:Npn \book_sidecite_separator:
  {
    \nobreak
    \hbox{\@textsuperscript{\normalfont\footnotesize\multfootsep}}
    \nobreak
  }
\bool_new:N \l_book_sidecite_first_bool
\cs_new_protected:Npn \book_sidecite_one:nnn #1#2#3
  {
    \book_sidecite_note:n
      {%
        \IfNoValueTF{#1}
          {\fullcite{#3}}
          {
            \IfNoValueTF{#2}
              {\fullcite[#1]{#3}}
              {\fullcite[#1][#2]{#3}}
          }
      }%
  }
\cs_new_protected:Npn \book_sidecite_each:nnn #1#2#3
  {
    \unskip\nobreak
    \bool_set_true:N \l_book_sidecite_first_bool
    \clist_map_inline:nn {#3}
      {
        \bool_if:NTF \l_book_sidecite_first_bool
          {\bool_set_false:N \l_book_sidecite_first_bool}
          {\book_sidecite_separator:}
        \book_sidecite_one:nnn {#1} {#2} {##1}
      }
  }
\cs_new_protected:Npn \book_scite_each:n #1
  {
    \unskip\nobreak
    \bool_set_true:N \l_book_sidecite_first_bool
    \clist_map_inline:nn {#1}
      {
        \bool_if:NTF \l_book_sidecite_first_bool
          {\bool_set_false:N \l_book_sidecite_first_bool}
          {\book_sidecite_separator:}
        \book_sidecite_note:n {\fullcite{##1}}
      }
  }
\NewDocumentCommand{\sidecite}{o o m}
  {\book_sidecite_each:nnn {#1} {#2} {#3}\ignorespaces}
\RenewDocumentCommand{\cite}{o o m}
  {\book_sidecite_each:nnn {#1} {#2} {#3}\ignorespaces}
% 句末引用把相邻标记和中文句号连续输出；源码中已有句号时将其吸收。
\NewDocumentCommand{\scite}{m t。}
  {
    \book_scite_each:n {#1}
    \nobreak。\ignorespaces
  }
\ExplSyntaxOff
\makeatother
